\chapter*{Résumé}

\noindent\textbf{Résumé}

\noindent Ce rapport présente un projet de fin d'études mené au sein de l'entité AI~\&~Data d'EY~Tunisie, portant sur la gestion du risque de crédit bancaire. Il développe une plateforme qui relie l'évaluation du risque de crédit, le calcul des pertes attendues selon IFRS~9, l'estimation du capital économique et l'analyse de la rentabilité, au service de deux usages complémentaires : l'évaluation d'une nouvelle demande de crédit (Origination) et le suivi d'un portefeuille déjà constitué (Portfolio Analytics). En l'absence d'historique réel de défaut, les modèles de probabilité de défaut et de perte en cas de défaut sont entraînés sur des historiques de crédit synthétiques, construits à partir d'une base réelle de comptes de dépôt ; les modèles de probabilité de défaut sont comparés selon une démarche champion/challenger propre à chaque parcours, tandis que la perte en cas de défaut est estimée par le modèle le mieux ajusté aux données parmi deux candidats. Ces estimations alimentent le calcul des pertes attendues, du capital et de la rentabilité, puis un mécanisme de décision explicable et gouverné, qui préserve l'intervention humaine. L'ensemble a été traduit en une architecture logicielle réelle, incluant un assistant conversationnel dédié à l'exploration des résultats, puis validé par une campagne de tests automatisés portant sur 1\,345 exécutions, dont 1\,325 réussies. Les limites méthodologiques et techniques identifiées, en particulier la nature synthétique des données et l'absence de déploiement en production, sont explicitées tout au long du rapport.

\medskip
\noindent\textbf{Mots-clés :} risque de crédit, probabilité de défaut, perte en cas de défaut, IFRS~9, capital économique, RAROC, aide à la décision, assistant conversationnel.

\bigskip
\noindent\textbf{Abstract}

\noindent This report presents an end-of-studies project carried out within the AI~\&~Data entity of EY~Tunisia, addressing bank credit risk management. It develops a platform that connects credit risk assessment, expected loss computation under IFRS~9, economic capital estimation and profitability analysis, serving two complementary uses: evaluating a new credit application (Origination) and monitoring an existing portfolio (Portfolio Analytics). In the absence of real default history, probability-of-default and loss-given-default models are trained on synthetic credit histories built from a real deposit-account database; probability-of-default models are compared under a champion/challenger approach specific to each pathway, while loss given default is estimated by the best-fitting model among two candidates. These estimates feed the computation of expected losses, capital and profitability, followed by an explainable, governed decision mechanism that preserves human intervention. The whole was translated into a real software architecture, including a conversational assistant dedicated to exploring the results, and validated through an automated test campaign covering 1,345 executions, of which 1,325 passed. The methodological and technical limitations identified, in particular the synthetic nature of the data and the absence of production deployment, are made explicit throughout the report.

\medskip
\noindent\textbf{Keywords:} credit risk, probability of default, loss given default, IFRS~9, economic capital, RAROC, decision support, conversational assistant.
